\documentclass[12pt,a4paper]{article}
\usepackage[UTF8]{ctex}
\usepackage{geometry}
\usepackage{graphicx}
\usepackage{float}
\usepackage{amsmath}
\usepackage{amsfonts}
\usepackage{amssymb}
\usepackage{booktabs}
\usepackage{multirow}
\usepackage{array}
\usepackage{tabularx}
\usepackage{longtable}
\usepackage{enumerate}
\usepackage{fancyhdr}
\usepackage{lastpage}
\usepackage{hyperref}
\usepackage[normalem]{ulem}

% 页面设置
\geometry{left=2.5cm,right=2.5cm,top=2.5cm,bottom=2.5cm}
\setlength{\headheight}{15pt}

% 页眉页脚设置
\pagestyle{fancy}
\fancyhf{}
\fancyhead[C]{实\hspace{0.5em}验\hspace{0.5em}原\hspace{0.5em}始\hspace{0.5em}记\hspace{0.5em}录}
\fancyfoot[L]{\parbox[t]{\textwidth}{实验人员：\underline{陈恪瑾}  学号：\underline{U202413925}  同组人员：\underline{黄伟峰}}}
\fancyfoot[R]{教师签字：\underline{\hspace{3cm}}}

\begin{document}

% 正文第一行标题
\noindent
{\Large \textbf{实验十一：}\underline{\textbf{串联谐振电路}}} \hfill {\Large \textbf{实验台编号：}\underline{\textbf{05}}}
\vspace{0.5cm}

\section*{一、实验注意事项}
\begin{enumerate}
    \item 测量时，应保证电源电压有效值不变。
    
    \item 电感线圈具有电阻，电路中$ R $包含了电感线圈的电阻。
    
    \item 采样电阻$ r $应远小于被测电路的入端阻抗。
    
\end{enumerate}

\section*{二、实验线路}
\begin{figure}[H]
    \centering
    \includegraphics[width=0.8\textwidth]{E11_preview/schematic.png}
    \caption{串联谐振实验电路}
    \label{fig:circuit}
\end{figure}

\section*{三、实验数据}
$ Q = 1.5, R = 510 $

\begin{longtable}{|c|*{9}{c|}}
\caption{幅频特性和相频特性数据} \label{tab:data} \\
\hline
 & $f_{l3}$ & $f_{l2}$ & $f_{lc}$ & $f_{l1}$ & $f_0$ & $f_{h1}$ & $f_{hc}$ & $f_{h2}$ & $f_{h3}$ \\
\hline
\endfirsthead

\multicolumn{10}{c}%
{\tablename\ \thetable\ -- 续表} \\
\hline
 & $f_{l3}$ & $f_{l2}$ & $f_{lc}$ & $f_{l1}$ & $f_0$ & $f_{h1}$ & $f_{hc}$ & $f_{h2}$ & $f_{h3}$ \\
\hline
\endhead

\hline
\multicolumn{10}{r}{续下页} \\
\endfoot

\hline
\endlastfoot

频率$f$ (Hz) &320 & 640& 854& 980& 1362& 1640& 1840& 2180& 3280\\
\hline
$u_{rpp}(V)$  & 9.87& 9.48& 9.05& 8.73& 7.71& 7.14& 6.71& 5.85& 4.31\\
\hline
$I(mA)$ & 2.72& 5.52& 7.75& 9.18& 10.96& 9.20& 7.74& 5.48& 2.75\\
\hline
$I/I_0$ & $\frac{1}{4}$& $\frac{1}{2}$& $\frac{1}{\sqrt{2}}$ & $\frac{1}{2^{\frac{1}{4}}}$& 1 & $\frac{1}{2^{\frac{1}{4}}}$& $\frac{1}{\sqrt{2}}$ & $\frac{1}{2}$& $\frac{1}{4}$\\
\hline
$\Phi_{u_s-i}(^\circ)$ & -80.12& -67.69& -54.0& -41.96& 0 & 25.88& 39.33& 53.68& 72.00\\
\hline
\end{longtable}


\end{document}
